\section{Three Processes and One Query}
\label{sec:ch07-three-processes}
\index{federated-wire sample|(}%

The sample is two ASP.NET Core projects and one container. Catalog listens on
port 5201 and owns \code{Product}. Reviews listens on 5202 and knows nothing
about products except their identifiers. The Cosmo Router listens on 3002 and
has never met either of them until it is handed a file.

Every capture in this chapter comes from those three at fixed versions:
HotChocolate 16.6.0 in both subgraphs, Cosmo Router 0.337.1 in the container,
and wgc 0.129.7 doing the composition. Router output in particular is
version-specific, and the plan of section~\ref{sec:ch07-the-plan} is a Cosmo
artefact rather than anything the Federation specification defines.

\index{subgraph!Catalog sample}%
Here is the part of Catalog's schema that is its own work. Everything the
federation package adds around it is section~\ref{sec:ch07-what-a-subgraph-owes}:

\begin{graphqlsdl}
type Query {
  products: [Product!]!
  productById(id: ID!): Product
}

type Product @key(fields: "id") {
  id: ID!
  title: String!
  price: Decimal!
}
\end{graphqlsdl}

\index{subgraph!Reviews sample}%
And here is Reviews:

\begin{graphqlsdl}
type Query {
  reviews: [Review!]!
}

type Product @key(fields: "id") {
  id: ID!
  reviews: [Review!]!
}

type Review {
  id: ID!
  rating: Int!
  body: String!
}
\end{graphqlsdl}

What matters here is the overlap. \code{Product} appears twice,
with the same key both times, and the two declarations share exactly one field
between them. Catalog has never heard of a review. Reviews has never heard of a
title.

\code{Review} carries no \code{@key}, and that is deliberate. It is a value type
in the sense chapter~\ref{ch:the-federation-model} gave the term: it lives in one
subgraph, no other service contributes to it, and the router has no way to ask
for one by identity. Section~\ref{sec:ch07-representations} shows what happens
when you ask anyway.

\subsection*{Composing them}

\index{composition!wgc router compose@\code{wgc router compose}}%
The router needs one file, and one command produces it:

\begin{minted}{text}
npx wgc router compose -i samples/federated-wire/graph.yaml \
                       -o samples/federated-wire/supergraph.json
\end{minted}

\code{graph.yaml} is the composer's input: a name, a routing URL and a schema
file for each subgraph.

\begin{minted}{yaml}
version: 1

subgraphs:
  - name: catalog
    routing_url: http://host.docker.internal:5201/graphql
    schema:
      file: ../../schema/samples/wire-catalog.graphql

  - name: reviews
    routing_url: http://host.docker.internal:5202/graphql
    schema:
      file: ../../schema/samples/wire-reviews.graphql
\end{minted}

\index{Cosmo!local composition}%
The Cosmo documentation is explicit that this command \enquote{does not interact
with the control plane and completely runs
locally}~\autocite{cosmo2026routercompose}. It needs no account and contacts no
registry. Two files in, one file out.

That output file is not a supergraph schema, which is what I expected the first
time. It is a \emph{router execution config}: eight kilobytes of JSON for these
two subgraphs, with the supergraph's SDL as one string inside it alongside the
datasource configuration the router's query planner needs.
Chapter~\ref{ch:composition} takes it apart. For now the only interesting thing
about it is the version it carries:

\begin{minted}{json}
"version": "00000000-0000-0000-0000-000000000000",
"compatibilityVersion": "1:0.63.2"
\end{minted}

All zeroes, because a locally composed config has no registry to get a version
from. The router logs that value at startup as its \code{config\_version}, which
is a small thing to know before you go looking for why every graph in your
development environment claims to be the same one.

\subsection*{Starting the router without an account}

\index{Cosmo!router configuration}%
The router runs from the repository's \code{docker-compose.yml}, behind a
compose profile so that nothing which starts Mosaic starts this. Three
environment variables:

\begin{minted}{yaml}
EXECUTION_CONFIG_FILE_PATH: "/etc/wire/supergraph.json"
DEV_MODE: "true"
LISTEN_ADDR: "0.0.0.0:3002"
\end{minted}

There is no \code{GRAPH\_API\_TOKEN}. The documentation says it \enquote{can be
empty when providing a static router configuration \dots\ but will disable the
default telemetry stack}~\autocite{cosmo2026routerconfig}, and the router says
the same thing rather more loudly on the way up:

\begin{minted}[fontsize=\footnotesize]{text}
WARN core/router.go:464 No graph token provided. The following Cosmo Cloud
features are disabled. Not recommended for Production.
{"features": ["Schema Usage Tracking", "Persistent operations",
              "Cosmo Cloud Tracing", "Cosmo Cloud Metrics"]}
WARN core/router.go:536 Development mode enabled. This should only be used for
testing purposes
INFO core/router.go:1722 Static execution config provided. Polling and watching
is disabled. Updating execution config is only possible by restarting the router
\end{minted}

\code{DEV\_MODE} is not decoration. It turns on Advanced Request Tracing without
an authentication check in front of it, and the query plans of
section~\ref{sec:ch07-the-plan} are part of that. Without it the header this
chapter leans on does nothing.

The third warning is the one that shapes the rest of the chapter. A statically
configured router neither polls nor watches, and it never asks the subgraphs
anything. Everything it knows about them, it read out of that file.

\subsection*{The query}

\index{federated query!first}%
With all three running, the client sends one document to one endpoint:

\begin{minted}{graphql}
{
  products {
    title
    price
    reviews {
      rating
      body
    }
  }
}
\end{minted}

Neither service can answer that. Catalog has the titles and prices and no
reviews; Reviews has the reviews and cannot list a product. The response comes
back anyway:

\begin{minted}[fontsize=\footnotesize]{json}
{"data":{"products":[
  {"title":"Oak dining table","price":749.00,"reviews":[
     {"rating":5,"body":"Survived a house move."},
     {"rating":3,"body":"Arrived with a chipped leg."}]},
  {"title":"Linen armchair","price":429.00,"reviews":[
     {"rating":4,"body":"Firmer than it looks."}]},
  {"title":"Brass floor lamp","price":189.00,"reviews":[]}]}}
\end{minted}

That is the router's own output with line breaks added; it arrives as one line.

The third product has never been reviewed, and the empty list is worth noticing
now because it costs a round trip later. Reviews was asked about it and answered
\enquote{none}, which is a different thing from not being asked.

Everything from here is the mechanism behind that answer.

\index{federated-wire sample|)}%
